\subsection{L1: constructibility without local representing schemes}

We retain all the hypotheses of the literal criterion in scan 14,
including actual condition 7 and finite-type tests in condition 8.
The counterexample in Theorem~\ref{A:main-counterexample} shows that
these hypotheses do not ensure the existence of actual modular local
models.  Nevertheless they do control the locus of such models when
a finite-type test and a section have already been given.

For a locally finite-type $S$-scheme $T$ and $\eta\in F(T)$,
write $E(T,\eta)$ for the following \emph{Artinian formally
\'etale locus}.  A point $t$ belongs to $E(T,\eta)$ if every
diagram consisting of an Artinian local $S$-algebra $D$, a nilpotent
ideal $J\subset D$, a map $\Spec(D/J)\to T$ centered at $t$,
and an element of $F(D)$ restricting to the given image of $\eta$
has a unique lift $\Spec D\to T$.  Arbitrary residue-field
extensions are included in these tests.  This condition is preserved
by base change and is local for \'etale changes of the source and
base.  For the latter assertion, extend an Artinian test by a finite
faithfully flat \'etale algebra to lift the chosen residue point,
use formal \'etaleness of the \'etale change to extend that lift
through its nilpotent ideal, and descend the unique resulting map.
Here $M(T,\eta)$ continues to denote actual local modularity,
including the field-of-definition requirement.

\begin{lemma}[Formal lifting through an open fpqc refinement]
\label{A:fpqc-open-formal-descent}
\status{Proved}
Let $A\to A'$ be flat, let $X$ be an $A$-scheme, and let
$C\subset X_{A'}$ be an open subscheme for which $C\to X$ is
fpqc.  Let $F$ be an fpqc sheaf over $A$ and $h_X\to F$ a
morphism.  If $h_C\to F_{A'}$ is formally \'etale on all
nilpotent test schemes, then $h_X\to F$ has the same property.
\end{lemma}
\begin{proof}
Consider a nilpotent closed immersion $Z_0\hookrightarrow Z$,
a map $g_0:Z_0\to X$, and $\zeta\in F(Z)$ matching $g_0$
on $Z_0$.  In $(Z_0)_{A'}$ the inverse image of $C$ is an open
subscheme $V_0$.  Since nilpotent thickenings have the same
underlying topological spaces, there is a unique open
$V\subset Z_{A'}$ with $V\times_ZZ_0=V_0$.
The map $V_0\to Z_0$ is the base change of $C\to X$, hence
fpqc.  The map $V\to Z$ is flat, is surjective because $V_0$
is, and is quasi-compact: quasi-compactness of this open can be
tested after the nilpotent reduction.  Thus it is fpqc.
Formal \'etaleness gives a unique lift $V\to C$ of $V_0\to C$
with image $\zeta|_V$.

On $V\times_ZV$ the two resulting maps to $X$ agree modulo
the nilpotent ideal.  Pair \emph{both} maps with the $A'$-structure
from the first factor.  Their reductions agree and land in the
same open $C\subset X_{A'}$; therefore both maps themselves
factor through $C$.  Both have the prescribed image in $F_{A'}$.
Formal uniqueness for $C$ makes them equal.  Fpqc descent of
maps into the scheme $X$ gives $g:Z\to X$; separatedness of the
sheaf $F$ checks that its image is $\zeta$.  Any two proposed
lifts become maps to that same open $C$ over $V$ and are equal
there, so they agree.  This proves uniqueness as well.
\end{proof}

\begin{proposition}[The infinitesimal consequence of one actual modular point]
\label{A:actual-pro-generization}
\status{Proved}
Assume conditions 1--7.  If $T$ is locally of finite type over
$S$ and $x\in M(T,\eta)$, then every generization of $x$ in
$T$ belongs to $E(T,\eta)$.
\end{proposition}
\begin{proof}
Put $s=f(x)$ and $A=\mathcal O_{S,s}$.  Condition 5 and
Lemma~\ref{A:artinian-base} represent $F_{\kappa(s)}$ by a
disjoint union of finite local schemes.  Modularity is preserved
by base change: for a local test one first uses the original
local universal property and then combines the resulting map
with the new base map; the field-of-definition condition picks
the required point of the fiber product.  Thus the map of the
fiber $T_s$ to this representing scheme is a local-ring
isomorphism at $x$.  In particular $T\to S$ is quasi-finite
at $x$, and $\kappa(x)/\kappa(s)$ is finite.

Let $\widehat A$ be the completion of $A$.  Over this henselian
local base, the quasi-finite branch of $T_{\widehat A}$ through
the point over $x$ has a finite local open neighborhood
$C=\Spec B$; this is the henselian finite-branch consequence of
Zariski's Main Theorem used in
Theorem~\ref{A:local-models-suffice} and
\cite[Lemma 10.153.4(3)]{StacksHenselianBranches}.
Its closed point is actual modular by the preceding base-change
argument.  Condition 7 therefore makes its section locally
modular at every point of $C$; the stated weaker \'etale variant
of 7 would also suffice here.

The morphism $h_C\to F_{\widehat A}$ is representable and
\'etale.  To check this rather than assume a representation
of $F$, for each noetherian $\widehat A$-scheme $Y$ with a
section of $F$ consider $h_C\times_Fh_Y$.  Conditions 1--4
are preserved by this fiber product, its Artinian nilpotent
lifting maps are bijective by the local modularity just proved,
and valuative uniqueness follows from the separated finite
scheme $C/\widehat A$.  Theorem~\ref{A:etale-corollary}
therefore represents it by a separated \'etale $Y$-scheme.
For arbitrary affine $Y$, finitarity descends its section to a
finite-type $\widehat A$-algebra, and base change gives the
same assertion; then use Zariski gluing.  In particular this
chart has formal \'etaleness on every nilpotent test.

Set $X=\Spec\mathcal O_{T,x}$.  The whole local scheme $C$
maps to $X$: every function invertible at $x$ becomes a unit
in the local ring $B$.  Thus $C$ is an open subscheme of
$X\times_A\widehat A$.  The map
$\mathcal O_{T,x}\to B$ is flat, since it is a localization
of the flat completion base change, and local, since its closed
point maps to $x$.  It is therefore faithfully flat; the map
of affine schemes $C\to X$ is fpqc.
Lemma~\ref{A:fpqc-open-formal-descent} makes $h_X\to F_A$
formally \'etale.  An Artinian local test centered at a
generization of $x$ factors through $X$, as does any nilpotent
lift of its map to $T$.  Its unique lifting property is therefore
the required one for $h_T\to F$.  This proves the assertion.
No claim that actual local modularity descends fpqc is used.
\end{proof}

\begin{lemma}[Reflection from one thickening along a stratum]
\label{A:generic-thickening-reflection}
\status{Proved}
Assume conditions 1--6 and 8.  Let $T$ be locally of finite type
over $S$, let $\eta\in F(T)$, and let $Z\subset T$ be an
irreducible closed subset with generic point $z$.
If $z\in E(T,\eta)$, there is a nonempty relatively open
subset of $Z$ contained in $E(T,\eta)$.
\end{lemma}
\begin{proof}
Write $s=f(z)$.  Over the Artinian base $\kappa(s)$ the functor
is actually represented by finite local schemes, by
Lemma~\ref{A:artinian-base}.  The definition of $E$ and the
infinitesimal criterion for \'etaleness show that the map of
the fiber $T_s$ to that representation is \'etale at $z$.
Consequently $T/S$ is quasi-finite at $z$.
After an \'etale neighborhood of $s$ and a restriction around
a point over $z$, isolate a finite branch
$T=\Spec C_R\to\Spec R$.  The finite branches and their open
inclusions descend from a strict henselization to a finite
\'etale stage, as in the finite-branch argument of
Theorem~\ref{A:local-models-suffice}; see also
\cite[Lemmas 10.153.4 and 10.153.6]{StacksHenselianBranches}.

We can also arrange that at $z$ the residue field of the source
equals the minimal field of definition of its image in $F$.
Indeed their extension is finite separable by \'etaleness on
the Artinian fiber.  Over a separable closure of the base
residue field both finite residue extensions are purely
inseparable, so that separable extension becomes trivial on
the chosen local factor.  This finite residue splitting already
occurs after a finite separable base extension, which is realized
by a finite-presentation \'etale neighborhood.  Restrict to the
corresponding finite branch.  These \'etale modifications have
open images and preserve and reflect the Artinian lifting
condition, so a relative open proved after them suffices.

Let $I\subset R$ be the prime defining the closure of the image
of $z$.  Choose $N$ from Lemma~\ref{A:uniform-AR} for the finite
$R$-module $C_R$.  Further flat \'etale refinements do not alter
the validity of this choice.  Over
$\Spec((R/I^N)_I)$ the base is Artinian.  Thus $F$ is actually
represented there.  The branch of $T$ at $z$ maps \'etale to
the finite local representing branch and has the same residue
field, so the map of local Artinian rings is an isomorphism.
This is actual generic modularity after the base change to
$R/I^N$, even though the original hypothesis at $z$ was only
the Artinian lifting condition.

Remove from $T\times_R\Spec(R/I^N)$ the other irreducible
components, retaining an irreducible open containing $z$.
Condition 8, with this finite-type test and the irreducible
noetherian base $\Spec(R/I^N)$, gives a nonempty open $W$ of
this component on which the section is actually modular, or
\'etale in the allowed weaker conclusion, over $R/I^N$.
We claim every $w\in W$ belongs to the original $E(T,\eta)$.

Put $q=f(w)$ and $R'=\widehat{R_q}$.  The finite algebra
$C_R\otimes_RR'$ is a product of finite local algebras.  Its
factor $C_w$ through $w$ is the completion of
$\mathcal O_{T,w}$ and is a direct summand of the fixed target
module $C_R\otimes_RR'$.
Conditions 4 and 5 give a finite complete local $R'$-algebra
$D_w$ representing the local Artinian deformation functor of
the field-valued point selected by $\eta(w)$.
For precision, take the compatible finite local representing
algebras over $R'/\mathfrak m_{R'}^n$.  A finite set of
generators at level one lifts through the surjections, and
successive correction of coefficients in powers of
$\mathfrak m_{R'}$ proves that their limit is finite over
the complete ring $R'$.  Reduction of the limit gives each
finite-level algebra.  This is the same finite inverse-limit
argument used in Theorem~\ref{A:pro7-representability}; it
does not assert actual modularity of the limit.
The section on $C_w$ induces a local map
\[
                         D_w\longrightarrow C_w.
\]

If the conclusion of 8 is actual modularity, its reduction
modulo $I^N$ identifies these two complete local prorepresenting
algebras, so
$D_w/I^ND_w\simeq C_w/I^NC_w$.
Apply Lemma~\ref{A:uniform-AR} over the flat noetherian local
map $R\to R'$, including its direct-summand assertion.
It gives $D_w\simeq C_w$.

The weaker \'etale conclusion is sufficient too.  In that
case the map of complete local deformation rings
$D_w/I^ND_w\to C_w/I^NC_w$ has unique lifting on local
Artinian tests.  Its residue-field extension is finite separable:
this may already be checked on the represented finite Artinian
fiber over the residue field of $R'$.
Let $D'_w$ be the
finite \'etale local extension of the complete henselian ring
$D_w$ with that residue field, supplied by
\cite[Lemma 10.153.7]{StacksFiniteEtaleHenselL1}.
Henselian lifting into $C_w$
gives a map $D'_w\to C_w$ extending $D_w\to C_w$ and the
chosen residue map.  Modulo $I^N$ it is an isomorphism.  Indeed,
both complete local algebras have the unique Artinian lifting
property with the prescribed residue embedding over
$D_w/I^ND_w$.  Apply the two universal properties to their
Artinian quotients, pass to the inverse limit, and use uniqueness
on every quotient to make the two maps inverse.  This argument
only uses tests centered at the closed point; it does not
assume a global Artinian lifting property at other primes.
Uniform Artin--Rees, now with finite source $D'_w$, gives
$D'_w\simeq C_w$.  Thus $D_w\to C_w$ is \'etale.

Either conclusion is exactly the original Artinian formal
\'etaleness at $w$.  Indeed every local Artinian $R_q$-test
factors uniquely through $\widehat{R_q}$, and every local
Artinian test map centered at $w$ factors uniquely through
$\widehat{\mathcal O_{T,w}}=C_w$.  Hence no descent of
actual modularity from a completion is being asserted.

For a finite morphism, primes over $I$ are incomparable.
Consequently, near $z$, the retained irreducible component of
the inverse image of $V(I)$ is precisely the original stratum
$Z$.  The open $W$ supplies the desired relative open after
the modifications, and its \'etale image supplies a nonempty
relative open of the original $Z$.  This proves the lemma.
\end{proof}

\begin{theorem}[L1, with the infinitesimal induction repaired]
\label{A:L1-general}
\status{Proved}
Let $S$ be locally noetherian and let $F$ satisfy conditions
1--8 of cote 131, scan 14, using the actual premise in 7 and
locally finite-type tests in 8.  For every locally finite-type
$S$-scheme $T$ and $\eta\in F(T)$, the actual modularity locus
$M(T,\eta)$ is open, and in particular locally constructible.
No actual representation of $F$ over the local bases is assumed.
The \'etale alternatives in the conclusions of 7 and 8 suffice.
\end{theorem}
\begin{proof}
Work on a noetherian affine piece of $T$ over an affine piece
of $S$.  Let $E=E(T,\eta)$ and let $U$ be the union of the
open subsets of $T$ contained in $E$.  Then $U$ is open and
every point of $U$ lies in $E$.
We first prove
\[
                              M(T,\eta)\subset U.           \tag{A.L1}
\]
If $x\in M(T,\eta)\setminus U$, choose an irreducible
component $Z$ of the closed complement $T\setminus U$
containing $x$, and let $z$ be its generic point.
Proposition~\ref{A:actual-pro-generization} gives $z\in E$.
Lemma~\ref{A:generic-thickening-reflection} supplies a
nonempty relative open $W\subset Z$ contained in $E$.
It contains $z$.  Delete the other components of the closed
complement and then shrink around $z$ so that its intersection
with the complement lies in $W$.
The resulting open neighborhood of $z$ has all its points in
$E$: its points are either in $U$ or in $W$.
It is therefore contained in $U$ by the definition of $U$,
contradicting $z\notin U$.  This proves (A.L1).

On $U$ the chart $h_U\to F$ is representable and \'etale.
Indeed, for a noetherian test $(Y,\zeta)$ its fiber product
$h_U\times_Fh_Y$ has conditions 1--4, has bijective Artinian
nilpotent lifting maps by the definition of $E$, and has
valuative uniqueness since $U/S$ is separated on the affine
piece under consideration.  Theorem~\ref{A:etale-corollary}
represents this fiber product by a separated \'etale scheme.
Finitary approximation and Zariski gluing extend the assertion
to arbitrary test schemes, exactly as in
Proposition~\ref{A:actual-pro-generization}.

The image $Q$ of this chart is an open subfunctor of $F$:
on each test scheme it is the open image of the represented
\'etale morphism.  The represented \'etale equivalence
relation $U\times_FU\rightrightarrows U$ makes $Q$ an
algebraic space \cite[Theorem 65.10.5]{Stacks02WW}.
It is locally quasi-finite over $S$ because
every point of $U$ lies in $E$, and condition 5, applied on
the residue-field fiber as above, makes $U/S$ locally
quasi-finite.

We spell out the separation needed to obtain a scheme.
For affine chart pieces $U_i,U_j$, the relation
$U_i\times_QU_j\to U_i$ is separated and \'etale; its map
to $U_i\times_SU_j$ is a monomorphism, and $U_j/S$ is affine
on the chosen pieces.  Its fibers are finite, since their
points inject into a fiber of the finite-type quasi-finite
scheme $U_j/S$.  Lemma~\ref{A:etale-finitefibres-qc} therefore
makes the relation quasi-compact.  Thus $Q/S$ is quasi-separated.
Condition 6 and the noetherian valuative criterion make $Q/S$
separated, as in the already proved separation step of
Theorem~\ref{A:local-models-suffice}.  The separated locally
quasi-finite space $Q$ is a scheme by
\cite[Proposition 67.50.2]{Stacks0417}.

Since $Q\subset F$ is open, local modularity of the section
at a point of $U$ is equivalent to an isomorphism of the local
rings of $U\to Q$.  An isomorphism at a point for this
locally finitely presented morphism spreads to a local
isomorphism on an open neighborhood, as in
Proposition~\ref{A:8-literal}.  Hence $M(T,\eta)\cap U$ is
open in $U$.  By (A.L1) this is all of $M(T,\eta)$, which
is therefore open in $T$.  The assertion is local on the
original schemes, so the affine reductions lose no points.
\end{proof}

The induction here propagates the Artinian lifting condition
on a fixed finite test module before constructing the actual
open chart.  It never identifies modularity after a closed
base change with the original modularity locus.  Thus it
repairs the specific transverse-infinitesimal error in the
naive noetherian-induction route stated in L1.
